% ------------------------------------------------------------------------------
% Formatvorlage für RFCCC1
% ------------------------------------------------------------------------------
%   erstellt von Christian Schiffler, 25.11.2010
%   basierend auf Arbeiten von Stefan Macke, 24.04.2009
%   http://blog.stefan-macke.de

% Dokumentenkopf ---------------------------------------------------------------
%   Diese Vorlage basiert auf "scrbook" aus dem koma-script.
% ------------------------------------------------------------------------------
\documentclass[
    fontsize=11pt, % Schriftgröße
    DIV10,
    a4paper, % Papierformat
    %oneside, % einseitiges Dokument
    openany, % remove empty pages after a chapter if chapter starts on the "left" hand
    titlepage, % es wird eine Titelseite verwendet
    parskip=half, % Abstand zwischen Absätzen (halbe Zeile)
    headings=normal, % Größe der Überschriften verkleinern
    listof=totoc, % Verzeichnisse im Inhaltsverzeichnis aufführen
    bibliography=totoc, % Literaturverzeichnis im Inhaltsverzeichnis aufführen
    index=totoc, % Index im Inhaltsverzeichnis aufführen
    captions=tableheading, % Beschriftung von Tabellen unterhalb ausgeben
    final % Status des Dokuments (final/draft)
]{scrbook}

% Umlaute ----------------------------------------------------------------------
%   Umlaute/Sonderzeichen wie äüöß direkt im Quelltext verwenden (CodePage).
%   Erlaubt automatische Trennung von Worten mit Umlauten.
% ------------------------------------------------------------------------------
\usepackage[utf8x]{inputenc}
\usepackage[T1]{fontenc}
\usepackage{textcomp} % Euro-Zeichen etc.

% Meta-Informationen -----------------------------------------------------------
%   Informationen über das Dokument, wie z.B. Titel, Autor, etc.
%   werden in der Datei Meta.tex definiert und können danach global
%   verwendet werden.
% ------------------------------------------------------------------------------
% Meta-Informationen -----------------------------------------------------------
%   Definition von globalen Parametern, die im gesamten Dokument verwendet
%   werden können (z.B auf dem Deckblatt etc.).
%
%   ACHTUNG: Wenn die Texte Umlaute oder ein Esszet enthalten, muss der folgende
%            Befehl bereits an dieser Stelle aktiviert werden:
%            \usepackage[latin1]{inputenc}
% ------------------------------------------------------------------------------
\newcommand{\titel}{Request for Contao Coder Comments 1 - RFCCC1}
\newcommand{\untertitel}{ Contao Extension Repository 3}
\newcommand{\autor}{Tristan Lins, Christian Schiffler}
\newcommand{\jahr}{2011}
\newcommand{\logo}{images/base/Title}
\newcommand{\deckblattlogo}{images/base/Title}


% benötigte Packages -----------------------------------------------------------
%   LaTeX-Packages, die benötigt werden, sind in die Datei Packages.tex
%   "ausgelagert", um diese Vorlage möglichst übersichtlich zu halten.
% ------------------------------------------------------------------------------
\chapter{Packages}

\section[sec:versioning]{Versioning}
The versioning is based on the classic version numbering scheme with extensions.

A version number is combined like the following.

\begin{lstlisting}[caption=Version number in EBNF]
DOT            = <the dot (".") character>
UPALPHA        = <any US-ASCII uppercase letter "A".."Z">
LOALPHA        = <any US-ASCII lowercase letter "a".."z">
DIGIT          = <any US-ASCII digit "0".."9">
ALPHA          = UPALPHA | LOALPHA
ALNUM          = ALPHA | DIGIT
MAJOR          = DIGIT
MINOR          = DIGIT
REVISION       = DIGIT
EXTENSION      = <"-"> *ALPHA *DIGIT
BUILD          = *DIGIT
HEADBUILD      = <"head"> [DOT *ALNUM]
REGULARBUILD   = <"alpha" DIGIT | "beta" DIGIT | "rc" DIGIT | "stable"> [DOT BUILD]
RELEASE        = HEADBUILD | REGULARBUILD

VERSION = MAJOR DOT MINOR [DOT REVISION] [EXTENSION] DOT RELEASE
\end{lstlisting}

The only mandatory fields in the version are major, minor and the release. 
The release number is combined of either the word "head" in conjunction with an alpha numeric revision suffix which refers to unstable releases or one of the words "alpha", "beta", "rc" or "stable" suffixed with one digit. The latter one can also be accompanied with a build number.
The revision number (also known as "maintenance" or "patchlevel") is entirely optional, as it is not used in most cases.

The optional extension component ist mostly useful for personal repositories, which serve altered versions of certain extensions. This is known from derivated works like Ubunutu is releasing altered debian packages with a suffix "ubuntu1" or similar.

Developers are encouraged to provides adjustments to upstream maintainers in favor of providing own versions of packages though.

\subsection[sec:versioning real word examples]{A real world versioning example}

Let's imagine an agency has altered an extension for ACME Inc. and is providing this extension via their own repository.
To eliminate problems with the vanilla extension in the main repository, the agency appends the extension "acme" to the version and results in the version
\DisplFile{1.2.3-acme1}.

Using this scheme one can prevent unwanted upgrading (see \ref{sec:versioncompatibility}). In contrast to the other components of the version, the build number is the only one that can not be influenced by the developer. The build number is incremented by one for every package being uploaded into the repository but not when the description or the properties are changing. This is in contrast to the old ER, which incremented the build number upon every change to the extension.

\section[sec:versioncompatibility]{Version compatibility}
Thanks to the \techTerm{extension} component of the versuion, it is possible to tag modified variants of an extension. These variants are in most cases incompatible to newer (vanilla) releases of the same extension.

To address these potentional problems, the following version compatibility rules are in place.

Versions are compatible if, and only if the following conditions are met:
\begin{itemize}
  \item ...the alphanumeric representation of the \techTerm{extension} component is identical.
  \item ...the extension component is omitted on both versions.
\end{itemize}

\small
\begin{longtable}{|p{0.25\textwidth}|p{0.05\textwidth}|p{0.25\textwidth}|}
\caption{Version compatibility examples (compatible: "="; incompatible: !)} \\
\hline
\label{tab:versioncompatexamples}
\textbf{Version 1} & \textbf{=/!} & \textbf{Version 2} \\
\hline
 1.2.3 & = & 4.5.6
\\ \hline
 1.2-extA.3 & = & 4.5-extA.6
\\ \hline
 1.2-extB.3 & ! & 4.5-extC.6
\\ \hline
 1.2-extD3.4 & = & 5.6-extD7.8
\\ \hline
 1.2-extE3.4 & ! & 5.6-extF7.8
\\ \hline
\end{longtable}
\normalsize

Please note that this only refers to the compatibility in respect to the version and that the compatibility can be further restricted by the extension meta file as described in \ref{sec:package information file}.

\section[sec:version comparision]{Version comparision}

When comparing compatible versions (according to the rules defined in \ref{sec:versioncompatibility}), one has to compare versions with extension component with a different approach than those without extension component.

\subsection[sec:version comparision without extension]{Version comparision without extension component}
A version without extension component will be matched according to the usual mechanism.

A given version A is bigger (aka. newer) than version B if, and only if:
\begin{itemize}
\item A.major>B.major
\item A.major=B.major \& A.minor>B.minor
\item A.major=B.major \& A.minor=B.minor \& A.revision>B.revision
\item A.major=B.major \& A.minor=B.minor \& A.revision=B.revision \& A.build>B.build
\end{itemize}

A given version A ist identical to version B if and only if:
\begin{itemize}
\item A.major=B.major \& A.minor=B.minor \& A.revision=B.revision \& A.build=B.build
\end{itemize}

\subsection[sec:version comparision with extension]{Version comparision with extension component}

If an \techTerm{extension} component is present, the following additional rules are to be applied in addition to the rules mentioned in \ref{sec:version comparision without extension}.

A given version A is bigger (aka. newer) than version B if, and only if:
\begin{itemize}
\item A.extensionA = B.extensionA \& A.extensionN > B.extensionN
\end{itemize}

A given version A ist identical to version B if and only if:
\begin{itemize}
\item A.extensionA = B.extensionA \& A.extensionN = B.extensionN
\end{itemize}

Hereby one must pay attention to the fact that the alphabetical representation and the numeric representation has to be evaluated seperately.

If the alphabethical representation (extensionA) is different for two versions, they have to be regarded as incompatible according to \ref{sec:versioncompatibility}. When no numeric representation (extensionN) is present for in a version with extension, the numeric part is implicit "1".

\small
\begin{longtable}{|p{0.25\textwidth}|p{0.25\textwidth}|p{0.25\textwidth}|}
\caption{Version comparision example} \\
\hline
\label{tab:versioncompareexamples}
\textbf{field} & \textbf{Version 1} & \textbf{Version 2} \\
\hline
full version & 1.2-acme3.4 & 1.2-acme4.1
\\ \hline
major & 1 & 1
\\ \hline
minor & 2 & 2
\\ \hline
extension & acme3 & acme4
\\ \hline
extensionA (alpha) & acme & acme
\\ \hline
extensionN (numeric) & 3 & 4
\\ \hline
\end{longtable}
\normalsize

\section[sec:packet structure]{Package structure}
Extensions are shipped in a package file. Such a package contains meta information, optional package managemen routines for installation, uninstallation and handling upgrades in addition to the application data.

The package file is a zip file with a changed extension to mark it as a package file. This approach is already known \ie from java (jar files).

For the file extension currently the following are considered (on the right hand are the file formats listed that are known to use this file extension):
\\
\begin{tabular}{lp{0.9\textwidth}}\\
.cext & none known so far\\
.cex  & INMOS Transputer Development System Occam User Program, ThumbsPlus File, The Currency Exchanger Rate File, CLAN Output File (Child Language Data Exchange System)\\
.ctp  & American Greetings CreataCard (Broderbund), BestAddress HTML Editor Combined Template File (Multimedia Australia Pty. Ltd.)\\
.ctx  & CTX Chemical File (Gasteiger Group), Chinese Character Input File, Compressed Text, GE Industrial Systems CIMPLICITY Text Version HMI Screen, Alphacam Compiled Text (Planit), Online Course Text (Microsoft Corporation), Pretty Good Privacy (PGP) Ciphertext File (PGP Corporation), Windows Terminal Server INI Backup File (Microsoft Corporation), Visual Basic User Control Binary File, CTRAN/W DEFINE Compressed Data File (Geo-Slope International), Nokia PC Suite Backup Contact File (Nokia)\\
\end{tabular}

A package contains the following files:
\begin{itemize}
\item PACKAGE - package information (required, see \ref{sec:package information file})
\item FILES - list of files contained in the package (required for certification, see \ref{sec:files information file})
\item LICENSE files - the license files for the extension (required according to \ref{sec:package information file})
\item ... any files the developer might need below the virtual root "CONTAO" which refers to the contao installation root.
\item ... any files the developer might need below the virtual root "DATA" for the user data root.
\end{itemize}

\subsection[sec:package information file]{PACKAGE information file}

The XML format is described in the schema file extension.xsd\supref{lst:extension.xsd}.

As usage example serves extension.xml\supref{lst:extension.xml}.

\subsubsection[sec:dependencies conflicts suggestions]{Package relationships - dependencies, conflicts and preferations}

A relationship of an extension to another extension can be defined in form of a \techTerm{dependency}, a \techTerm{conflict} or a \techTerm{preferation}.
Each of these relationships can be towards any version of the mentioned extension or even be narrowed down to certain version numbers.

\subsubsection[sec:versioning of relationships]{Versioning of relationships}

As dependencies and conflicts are often regarding certain versions or a certain range of versions.

In the definition of an relationship, we may use the following attributes:
\begin{itemize}
\item \techTerm{match} which means an exact match on the version (X=Y).
\item \techTerm{lowest} which is implicit inclusive and tells the lower end of the range (X>=Y).
\item \techTerm{highest} which is implicit exclusive and tells the upper end of the range (X<Y).
\end{itemize}
The attributes \techTerm{lowest} and \techTerm{highest} do not have to be used in combination, therefore it is possible to define an open range.
When used in combination, they define a closed range.

The attributes \techTerm{lowest} and \techTerm{highest} must not be used in conjunction with the attribute \techTerm{match}. If used in conjunction with \techTerm{match}, the exact match takes precedence over the other attributes and therefore is always exact.

In the package information file, this might look like the following when defining a conflict for a certain db driver (in this case mysql).
\lstset{language=XML} 
\begin{lstlisting}[caption=XML example: conflict example mysql]
...
<ctx:conflict name="db:mysql">
    <ctx:version match="4"/>
    <ctx:version lowest="5" highest="5.1"/>
    <ctx:version lowest="5.2"/>
</ctx:conflict>
...
\end{lstlisting}

In above example definition we have a conflict for:
\begin{itemize}
\item all version 4 drivers.
\item all versions 5 to 5.1 ( 5>=X<5.1)
\item all versions from 5.2 on upwards (5.2>=X)
\end{itemize}
With above condition we define a conflict for all versions of MySQL except for version 5.1 (we could have done it more simple though).

\subsubsection[sec:handling of versions with extension component]{Handling of versions with extension component}

In general the extension component of a version is ignored as it makes common sense that an altered version of an extension is compatible to it's unaltered counterpart from an external point of view. This behaviour can be altered by defining the extension version in the relationship explicitely. If specified, the match for the version must be performed including the extension part.

\subsubsection[sec:special relationships]{Special relationships}

Package relationships can not only refer other packages but some other namespaces as well.
This is useful to restrict dependency on some php features, a certain database driver, some PEAR package or something else that has been defined in the specification.

Currently the following namespaces exist.\\
\begin{tabular}{ll}\\
contao         & The contao core.\\
php:[feature]  & Relates to a certain php feature. See php features in \ref{sec:package relationship namespaces php features}\\
pear:[package] & Relates to a certain pear package. See pear packages in \ref{sec:package relationship namespaces pear packages}.\\
db:type        & Relates to a certain database, \ie db:mysql. See database drivers in \ref{sec:package relationship namespaces database drivers}.\\
\end{tabular}

\subsubsection[sec:alternative dependencies]{Alternative dependencies}
Often it is required to have a certain feature available but multiple ways exist to get this feature available when depending on some PHP feature but to also provide a wrapper class as workaround if it is not available.

Instead of defining features within the repository client and therefore having to update the client all the time when a new feature has been defined, we provide the developer with the possibility to define alternative pathes by specifying groups of dependencies.

When defining groups of dependencies, the first positive match found is accepted and the group evaluation will be stopped. Note that the order in the XML has no meaning and therefore the attribute \techTerm{preferation} must be used to define which dependency is the most preferable.

\lstset{language=XML} 
\begin{lstlisting}[caption=XML example: alternative relationships]
...
<ctx:group preferation="php:soap nusoap">
<ctx:dependency name="php:soap">
    <ctx:description>I love soap, thats because i need this php extension.</ctx:description>
</ctx:dependency>
<ctx:dependency name="nusoap">
    <ctx:description>If you dont have soap, I pleased with an alternative implementation.</ctx:description>
</ctx:dependency>
</ctx:group>
...
\end{lstlisting}
This example describes, that it depends on the PHP SOAP extension but alternatively can also use the extension "nusoap" if the first is not available.

\subsubsection[sec:upgrade and downgrade compatibility]{Upgrade and downgrade compatibility}
In many cases upgrades and downgrades may cause problems. This is even more the case, when the data schema is changing.
That's the reason why the developer needs a way to define what potentional threat an upgrade or downgrade of a certain package implies.

This potentional threat can be categorized in the following sections:

\begin{tabular}{llp{0.55\textwidth}}\\
\textbf{Category} & \textbf{threat factor} & \textbf{Description}\\
lossless  & none & Up- and downgrade can be performed without data loss.\\
dropUnused  & almost none & Up- and downgrade will remove data that will not be in use anymore.\\
partial  & might be dangerous & Up- and downgrade might remove user data. Therefore manual intervention is suggested.\\
undefined & dragons live here! &This setting is implicit defined when the developer has not defined a setting at all. The result might be user data loss, in worst case scenario resulting in total data loss.\\
\end{tabular}

Whenever a developer has not defined a compatibility, the following automatic assumptions are made:\\
When upgrading, the compatibility is defaulting to "partial".
When downgrading, the compatibility is defaulting to "undefined".
The extension repository manager will then display a message where the up/downgrade compatibility is shown in addition to the description text of the compatibility.

As "undefined" is an implicit type, the developer can not add a description text. For "undefined" compatibility the extension repository client will provide two possibilities:
\begin{enumerate}
\item soft-update - upgrading without removing any user data (hence by updating the schema still some data might get lost).
\item hard-update - remove the old version including all user data and reinstall the other version then (possibly the worst attempt).
\end{enumerate}

A corresponding upgrade/downgrade compatibility might look like the following:
\lstset{language=XML} 
\begin{lstlisting}[caption=XML example: upgrade/downgrade compatibility]
...
<ctx:extension version="1.2.3">
    <ctx:upgrade version="1.1" userdata="lossless"/>
    <ctx:upgrade version="1.0" userdata="partial"/>
    <ctx:downgrade version="1.2" userdata="lossless"/>
    <ctx:downgrade version="1.1" userdata="dropUnused"/>
    <ctx:downgrade version="1.0" userdata="partial"/>
</ctx:extension>
...
\end{lstlisting}

\todo{translation}

Dabei beschreibt die Version immer die jüngste (Upgrade) oder älteste (Downgrade) Version, bis zu der diese Regel gilt. In dem Beispiel hat die Erweiterung die Version 1.2.3.
Die Upgrade Fähigkeit besagt:
Von allen Versionen 1.1 < 1.2.3 kann ohne Datenverlust aktualisiert werden.
Von allen Versionen 1.0 < 1.1 kann mit Nutzdatenverlust aktualisiert werden.
Die Downgrade Fähigkeit besagt:
Zu allen Versionen 1.2 < 1.2.3 kann ohne Datenverlust aktualisiert werden.
Zu allen Versionen 1.1 < 1.2 kann mit dem Verlust ungenutzer Daten aktualisiert werden.
Zu allen Versionen 1.0 < 1.1 kann mit Nutzdatenverlust aktualisiert werden.
Die Reihenfolge ist demnach egal, die Upgrade Definitionen werden zuerst nach Version absteigend sortiert und bilden dann immer zur nächsten Definition einen Versionsbereich, für den die Regel gilt.

\subsection[sec:files information file]{FILES - package contents information file}
\todo{translation}
Die Paketinhaltsinformation beinhaltet eine Liste aller Paketdateien (mit Ausnahme der FILES Datei selbst) mit samt Prüfsumme. Die Datei dient der Validierung der Installation. Sie ist erforderlich, wenn ein signiertes Paket erzeugt wird, um mit dem Extension Manager eine zuverlässige Integritätsprüfung zu gewährleisten. Bei einem unsigniertem Paket kann die Datei mitgeliefert werden, ist aber nicht erforderlich. Wird ein unsigniertes Paket installiert, erstellt der Paket Manager die Datei bei der Installation.
Beispieldatei

\lstset{language=XML} 
\begin{lstlisting}[caption=FILES example: package contents information file]
...
4543a0fe56f27277948469c15284ac25 CONTAO/system/modules/myExtension/Class1.php
a692ffb1c6026092d916872b2114941c CONTAO/system/modules/myExtension/Class2.php
2709656afab36a8f635206c6956cbf53 CONTAO/system/modules/myExtension/templates/class1\_template.tpl
3fced4c880f491dce98c116266361f1c CONTAO/system/modules/myExtension/templates/class2\_template.tpl
b52f2d57d10c4f7ee67a7eb9615d5d24 LGPL
2e7781b7b08ab489e08526f43adf5b0c PACKAGE
...
\end{lstlisting}





\section[sec:package relationship namespaces]{Package relationship namespaces}

\subsection[sec:package relationship namespaces php features]{PHP features}

\todo{this section is currently undocumented}

\subsection[sec:package relationship namespaces pear packages]{PEAR features}

\todo{this section is currently undocumented}

\subsection[sec:package relationship namespaces database drivers]{Database driver}

\todo{this section is currently undocumented}


% Erstellung eines Index und Abkürzungsverzeichnisses aktivieren ---------------
\makeindex
\makenomenclature
\makeglossaries

% Kopf- und Fußzeilen, Seitenränder etc. ---------------------------------------
% Zeilenabstand 1,5 Zeilen -----------------------------------------------------
%\onehalfspacing
\spacing{1.0}

% Seitenränder -----------------------------------------------------------------
\setlength{\topskip}{\ht\strutbox} % behebt Warnung von geometry
\geometry{paper=a4paper}

% Kopf- und Fußzeilen ----------------------------------------------------------
\pagestyle{scrheadings}
% Kopf- und Fußzeile auch auf Kapitelanfangsseiten
\renewcommand*{\chapterpagestyle}{scrheadings} 
% Schriftform der Kopfzeile
\renewcommand{\headfont}{\normalfont}

% Kopfzeile
\ihead{\large{\textsc{\titel}}\\ \small{\untertitel} \\[2ex] \textit{\headmark}}
\chead{}
%\ohead{\includegraphics[scale=0.075]{\logo}}
\setlength{\headheight}{21mm} % Höhe der Kopfzeile
% Kopfzeile über den Text hinaus verbreitern
%\setheadwidth[0pt]{textwithmarginpar} 
\setheadsepline[text]{0.4pt} % Trennlinie unter Kopfzeile

% Fußzeile
\ifoot{\copyright\ \autor}
\cfoot{}
\ofoot{\pagemark}

% sonstige typographische Einstellungen ----------------------------------------

% erzeugt ein wenig mehr Platz hinter einem Punkt
\frenchspacing 

% Schusterjungen und Hurenkinder vermeiden
\clubpenalty = 10000
\widowpenalty = 10000 
\displaywidowpenalty = 10000

% Quellcode-Ausgabe formatieren
\lstset{numbers=left, numberstyle=\tiny, numbersep=5pt, breaklines=true}
\lstset{emph={square}, emphstyle=\color{red}, emph={[2]root,base}, emphstyle={[2]\color{blue}}}

% Fußnoten fortlaufend durchnummerieren
\counterwithout{footnote}{chapter}

%% ... und nimm alle 4 Ebenen in das Inhaltsverzeichnis auf.
\setcounter{tocdepth}{4}

% eigene Definitionen für Silbentrennung
\input{texbase/Hyphenation}

% eigene LaTeX-Befehle
% Eigene Befehle und typographische Auszeichnungen für diese

% einfaches Wechseln der Schrift, z.B.: \changefont{cmss}{sbc}{n}
\newcommand{\changefont}[3]{\fontfamily{#1} \fontseries{#2} \fontshape{#3} \selectfont}

% Abkürzungen mit korrektem Leerraum 
\newcommand{\ua}{\mbox{u.\,a.\ }}
\newcommand{\zB}{\mbox{z.\,B.\ }}
\newcommand{\dahe}{\mbox{d.\,h.\ }}
\newcommand{\Vgl}{Vgl.\ }
\newcommand{\bzw}{bzw.\ }
\newcommand{\evtl}{evtl.\ }
\newcommand{\ie}{\mbox{i.\,e.\ }}

\newcommand{\illustration}[1]{illustration~\ref{fig:#1}}

\newcommand{\bs}{$\backslash$}

\newcommand{\sectiontitle}{}
\let\TeXsection\section
\renewcommand{\section}[2][\empty]{\ifthenelse{\equal{#1}{\empty}}{\TeXsection{#2}\renewcommand{\sectiontitle}{#2}}{\TeXsection{#2}\label{#1}\renewcommand{\sectiontitle}{#1}}}

\newcommand{\subsectiontitle}{}
\let\TeXsubsection\subsection
\renewcommand{\subsection}[2][\empty]{\ifthenelse{\equal{#1}{\empty}}{\TeXsubsection{#2}\renewcommand{\subsectiontitle}{#2}}{\TeXsubsection{#2}\label{#1}\renewcommand{\subsectiontitle}{#1}}}

\let\TeXsubsubsection\subsubsection
\renewcommand{\subsubsection}[2][\empty]{\ifthenelse{\equal{#1}{\empty}}{\TeXsubsubsection{#2}}{\TeXsubsubsection{#2}\label{#1}}}

\newcommand{\supref}[1]{\textsuperscript{\ref{#1}}}

% erzeugt ein Listenelement mit fetter Überschrift 
\newcommand{\itemd}[2]{\item{\textbf{#1}}\\{#2}}

% einige Befehle zum Zitieren --------------------------------------------------
\newcommand{\Zitat}[2][\empty]{\ifthenelse{\equal{#1}{\empty}}{\citep{#2}}{\citep[#1]{#2}}}

% zum Ausgeben von Autoren
\newcommand{\AutorName}[1]{\textsc{#1}}
\newcommand{\Autor}[1]{\AutorName{\citeauthor{#1}}}

% verschiedene Befehle um Wörter semantisch auszuzeichnen ----------------------
\newcommand{\newDef}[1]{\textbf{#1}}
\newcommand{\techTerm}[1]{\textit{#1}}

\newcommand{\DisplInput}[1]{\texttt{#1}}
\newcommand{\DisplCode}[1]{\texttt{#1}}
\newcommand{\DisplFile}[1]{\texttt{#1}}

\newcommand{\Datentyp}[1]{\textsf{#1}}
\newcommand{\XMLElement}[1]{\textsf{#1}}
\newcommand{\Webservice}[1]{\textsf{#1}}

% Deklaration von Hinweistexten
\newcommand{\hint}[1]{
\begin{tabular}{w{1.5cm}p{12cm}}\\
\scalebox{0.125}{\includegraphics{images/base/Hint_Notification}} & #1
\end{tabular}
}

% Deklaration von Warnungen
\newcommand{\warning}[1]{
\begin{tabular}{w{1.5cm}p{12cm}}\\
\scalebox{0.125}{\includegraphics{images/base/Warning_Notification}} & #1
\end{tabular}
}

% einbinden eines emp Bildes.
\newcommand{\useempcenter}[2]{\begin{figure}[h!]\centering\scalebox{0.8}{\empuse{#1}}\caption{#2}\end{figure}}

\newcommand{\useempfloat}[3]{\begin{wrapfigure}{#1}{0.4\textwidth}\begin{center}\empuse{#2}\end{center}\caption{#3}\end{wrapfigure}}

\newcommand{\useemp}[3][\empty]{\ifthenelse{\equal{#1}{\empty}}{\useempcenter{#2}{#3}}{\useempfloat{#1}{#2}{#3}}}

% Wrapper für useemp für uml Klassendiagramme
\newcommand{\umlclass}[3][\empty]{\ifthenelse{\equal{#1}{\empty}}{\useemp{#2}{#3}}{\useemp[#1]{#2}{#3}}}

% Wrapper um listings einzubinden
\newcommand{\phplisting}[3]{\lstset{language=PHP, basicstyle=\footnotesize, showstringspaces=false, tabsize=2}\lstinputlisting[label=#1,caption=#2]{#3}}
\newcommand{\xmllisting}[3]{\lstset{language=XML, basicstyle=\footnotesize, showstringspaces=false, tabsize=2}\lstinputlisting[label=#1,caption=#2]{#3}}


% Das eigentliche Dokument -----------------------------------------------------
%   Der eigentliche Inhalt des Dokuments beginnt hier. Die einzelnen Seiten
%   und Kapitel werden in eigene Dateien ausgelagert und hier nur inkludiert.
% ------------------------------------------------------------------------------
\begin{document}

% auch subsubsection nummerieren
\setcounter{secnumdepth}{3}
\setcounter{tocdepth}{3}

% Deckblatt und Abstract ohne Seitenzahl
\ofoot{}
\thispagestyle{plain}
\begin{titlepage}

\begin{center}

\huge{\textbf{\titel}}\\[1.5ex]
\LARGE{\textbf{\untertitel}}\\[6ex]
%\includegraphics[scale=0.2]{\deckblattlogo}\\[6ex]

\normalsize
\autor\\[1.2ex]
\copyright\ \jahr\\[9ex]

\end{center}

\singlespacing
\small
\noindent 
\begin{center}
This work is licensed under a \\ \href{http://creativecommons.org/licenses/by-sa/3.0/}{Creative Commons Attribution-ShareAlike 3.0 Unported License}.
\end{center}

A human-readable summary of the Legal Code:

You are free:
\begin{itemize}
\item to Share — to copy, distribute and transmit the work
\item to Remix — to adapt the work
\item to make commercial use of the work
\end{itemize}

Under the following conditions:
\begin{itemize}
\item Attribution — You must attribute the work in the manner specified by the author or licensor (but not in any way that suggests that they endorse you or your use of the work).
\item Share Alike — If you alter, transform, or build upon this work, you may distribute the resulting work only under the same or similar license to this one.
\end{itemize}

Read the license for further details.
\href{http://creativecommons.org/licenses/by-sa/3.0/legalcode}{http://creativecommons.org/licenses/by-sa/3.0/legalcode}

\end{titlepage}

\chapter{Abstract}

This document serves as a discussion and documentation foundation for the new Extension Repository version 3 which shall replace the existing repository for Contao.

This document will be proposed for discussion at Contao Camp 2011.

The first goal of this document is, to identify problems with the current extension repository and provide solutions to them.

To fulfil this goal, the following conditions have to be met:
\begin{itemize}
\item Compatibility to Contao version 2 and version 3 (upcoming).
\item Be more flexible than the current system.
\item (almost) anyone shall be able to run his/her own repository.
\item Installation and uninstallation routines to aid extensions in setup.
\item Standardized package format to ease installation.
\item Improved dependencies.
\item Add possibility to define conflicting extensions.
\end{itemize}

The open discussion on this document is currently taking place in the 
\hyperlink{http://www.contao-community.de/showthread.php?22031-Gruppe-neuer-ER-Client-und-ER3}{German Contao Community}

\warning{This document is still a work in progress and should not be used for implementation yet as all parts are still subject to change. Once the document is finalized, it will get a version number and will be released along with some example implementation of the services.}
\ofoot{\pagemark}


% arabische Seitenzahlen im Hauptteil ------------------------------------------
\clearpage
\pagenumbering{arabic}
\tableofcontents % Inhaltsverzeichnis

% UML support
\begin{empfile}
\begin{empcmds}
input metauml;
\end{empcmds}
% die Inhaltskapitel werden in "Content.tex" inkludiert -------------------------
% Hier können die einzelnen Kapitel inkludiert werden. Sie müssen in den 
% entsprechenden .TEX-Dateien vorliegen. Die Dateinamen können natürlich 
% angepasst werden.

\chapter{Changelog}

\newcommand{\ChangeLogEntry}[3]{\subparagraph{Changes by #1 on #2} \begin{itemize} #3 \end{itemize} }

\tiny
\ChangeLogEntry{Tristan Lins}{08.09.2011}{
	\item Initial draft containing document description and package definition.
}
\ChangeLogEntry{Tristan Lins}{09.09.2011}{
	\item added PACKAGE XML information file.
	\item small adjustments in PACKAGE INI Format.
	\item elaborated usage of versions with an extension component.
}
\ChangeLogEntry{Tristan Lins}{10.09.2011}{
	\item Added first description of ER3 Repository, ER3 Repository Manager and ER3 Portal.
}
\ChangeLogEntry{Tristan Lins}{11.09.2011}{
	\item Removed INI Meta file specification.
	\item Reworked PACKAGE information file.
	\item Reworked XML schema.
	\item Reworked XML example.
}
\ChangeLogEntry{Tristan Lins}{12.09.2011}{
	\item Reworked XML Schema and example.
	\item added "versioning of dependencies" and "up/downgrade compatibility" to PACKAGE information.
	\item Removed description of license file and EULA.
}
\ChangeLogEntry{Christian Schiffler}{13.09.2011-14.09.2011}{
	\item Transition to \LaTeX{}, 
	\item translation to English and major rephrasing of all sections
	\item removed \techTerm{lowest\_bound} and \techTerm{highest\_bound} from \DisplFile{extension.xsd}.
}
\ChangeLogEntry{Tristan Lins}{15.09.2011}{
	\item Reworked Repository Portal\supref{sec:repository portal} section
}
\ChangeLogEntry{Christian Schiffler}{16.09.2011}{
	\item finished translation to English (mainly repository\supref{sec:repository} and repository manager \supref{sec:repository manager}).
}
\ChangeLogEntry{Tristan Lins}{19.09.2011}{
	\item Removed package scrhack, it generates the firstpage with "stylistings".
	\item Added first description of "Installation"\supref{sec:installation}, "Installation Process"\supref{sec:installation process} and "Installer"\supref{sec:installer_interfaces_and_classes}.
}
\ChangeLogEntry{Tristan Lins}{25.09.2011}{
	\item Rework complete chapter "Installation"
	\item Added section "ER3 Manager"\supref{sec:er3 manager}.
	\item Added section "Contao Installation Manager"\supref{sec:conatao installation manager}.
	\item Rename section "Installer" to "Installer interfaces and classes"\supref{sec:installer interfaces and classes}.
	\item Added subsection "InstallerInformation"\supref{sec:installerinformation}.
}
\ChangeLogEntry{Tristan Lins}{03.10.2011}{
	\item Added subsection "Extension Builder"\supref{sec:extension builder} to section "Repository Portal"\supref{sec:repository portal}
	\item Added section "Namespaces"\supref{sec:namespaces}.
}

\normalsize
\chapter{Packages}

\section[sec:versioning]{Versioning}
The versioning is based on the classic version numbering scheme with extensions.

A version number is combined like the following.

\begin{lstlisting}[caption=Version number in EBNF]
DOT            = <the dot (".") character>
UPALPHA        = <any US-ASCII uppercase letter "A".."Z">
LOALPHA        = <any US-ASCII lowercase letter "a".."z">
DIGIT          = <any US-ASCII digit "0".."9">
ALPHA          = UPALPHA | LOALPHA
ALNUM          = ALPHA | DIGIT
MAJOR          = DIGIT
MINOR          = DIGIT
REVISION       = DIGIT
EXTENSION      = <"-"> *ALPHA *DIGIT
BUILD          = *DIGIT
HEADBUILD      = <"head"> [DOT *ALNUM]
REGULARBUILD   = <"alpha" DIGIT | "beta" DIGIT | "rc" DIGIT | "stable"> [DOT BUILD]
RELEASE        = HEADBUILD | REGULARBUILD

VERSION = MAJOR DOT MINOR [DOT REVISION] [EXTENSION] DOT RELEASE
\end{lstlisting}

The only mandatory fields in the version are major, minor and the release. 
The release number is combined of either the word "head" in conjunction with an alpha numeric revision suffix which refers to unstable releases or one of the words "alpha", "beta", "rc" or "stable" suffixed with one digit. The latter one can also be accompanied with a build number.
The revision number (also known as "maintenance" or "patchlevel") is entirely optional, as it is not used in most cases.

The optional extension component ist mostly useful for personal repositories, which serve altered versions of certain extensions. This is known from derivated works like Ubunutu is releasing altered debian packages with a suffix "ubuntu1" or similar.

Developers are encouraged to provides adjustments to upstream maintainers in favor of providing own versions of packages though.

\subsection[sec:versioning real word examples]{A real world versioning example}

Let's imagine an agency has altered an extension for ACME Inc. and is providing this extension via their own repository.
To eliminate problems with the vanilla extension in the main repository, the agency appends the extension "acme" to the version and results in the version
\DisplFile{1.2.3-acme1}.

Using this scheme one can prevent unwanted upgrading (see \ref{sec:versioncompatibility}). In contrast to the other components of the version, the build number is the only one that can not be influenced by the developer. The build number is incremented by one for every package being uploaded into the repository but not when the description or the properties are changing. This is in contrast to the old ER, which incremented the build number upon every change to the extension.

\section[sec:versioncompatibility]{Version compatibility}
Thanks to the \techTerm{extension} component of the versuion, it is possible to tag modified variants of an extension. These variants are in most cases incompatible to newer (vanilla) releases of the same extension.

To address these potentional problems, the following version compatibility rules are in place.

Versions are compatible if, and only if the following conditions are met:
\begin{itemize}
  \item ...the alphanumeric representation of the \techTerm{extension} component is identical.
  \item ...the extension component is omitted on both versions.
\end{itemize}

\small
\begin{longtable}{|p{0.25\textwidth}|p{0.05\textwidth}|p{0.25\textwidth}|}
\caption{Version compatibility examples (compatible: "="; incompatible: !)} \\
\hline
\label{tab:versioncompatexamples}
\textbf{Version 1} & \textbf{=/!} & \textbf{Version 2} \\
\hline
 1.2.3 & = & 4.5.6
\\ \hline
 1.2-extA.3 & = & 4.5-extA.6
\\ \hline
 1.2-extB.3 & ! & 4.5-extC.6
\\ \hline
 1.2-extD3.4 & = & 5.6-extD7.8
\\ \hline
 1.2-extE3.4 & ! & 5.6-extF7.8
\\ \hline
\end{longtable}
\normalsize

Please note that this only refers to the compatibility in respect to the version and that the compatibility can be further restricted by the extension meta file as described in \ref{sec:package information file}.

\section[sec:version comparision]{Version comparision}

When comparing compatible versions (according to the rules defined in \ref{sec:versioncompatibility}), one has to compare versions with extension component with a different approach than those without extension component.

\subsection[sec:version comparision without extension]{Version comparision without extension component}
A version without extension component will be matched according to the usual mechanism.

A given version A is bigger (aka. newer) than version B if, and only if:
\begin{itemize}
\item A.major>B.major
\item A.major=B.major \& A.minor>B.minor
\item A.major=B.major \& A.minor=B.minor \& A.revision>B.revision
\item A.major=B.major \& A.minor=B.minor \& A.revision=B.revision \& A.build>B.build
\end{itemize}

A given version A ist identical to version B if and only if:
\begin{itemize}
\item A.major=B.major \& A.minor=B.minor \& A.revision=B.revision \& A.build=B.build
\end{itemize}

\subsection[sec:version comparision with extension]{Version comparision with extension component}

If an \techTerm{extension} component is present, the following additional rules are to be applied in addition to the rules mentioned in \ref{sec:version comparision without extension}.

A given version A is bigger (aka. newer) than version B if, and only if:
\begin{itemize}
\item A.extensionA = B.extensionA \& A.extensionN > B.extensionN
\end{itemize}

A given version A ist identical to version B if and only if:
\begin{itemize}
\item A.extensionA = B.extensionA \& A.extensionN = B.extensionN
\end{itemize}

Hereby one must pay attention to the fact that the alphabetical representation and the numeric representation has to be evaluated seperately.

If the alphabethical representation (extensionA) is different for two versions, they have to be regarded as incompatible according to \ref{sec:versioncompatibility}. When no numeric representation (extensionN) is present for in a version with extension, the numeric part is implicit "1".

\small
\begin{longtable}{|p{0.25\textwidth}|p{0.25\textwidth}|p{0.25\textwidth}|}
\caption{Version comparision example} \\
\hline
\label{tab:versioncompareexamples}
\textbf{field} & \textbf{Version 1} & \textbf{Version 2} \\
\hline
full version & 1.2-acme3.4 & 1.2-acme4.1
\\ \hline
major & 1 & 1
\\ \hline
minor & 2 & 2
\\ \hline
extension & acme3 & acme4
\\ \hline
extensionA (alpha) & acme & acme
\\ \hline
extensionN (numeric) & 3 & 4
\\ \hline
\end{longtable}
\normalsize

\section[sec:packet structure]{Package structure}
Extensions are shipped in a package file. Such a package contains meta information, optional package managemen routines for installation, uninstallation and handling upgrades in addition to the application data.

The package file is a zip file with a changed extension to mark it as a package file. This approach is already known \ie from java (jar files).

For the file extension currently the following are considered (on the right hand are the file formats listed that are known to use this file extension):
\\
\begin{tabular}{lp{0.9\textwidth}}\\
.cext & none known so far\\
.cex  & INMOS Transputer Development System Occam User Program, ThumbsPlus File, The Currency Exchanger Rate File, CLAN Output File (Child Language Data Exchange System)\\
.ctp  & American Greetings CreataCard (Broderbund), BestAddress HTML Editor Combined Template File (Multimedia Australia Pty. Ltd.)\\
.ctx  & CTX Chemical File (Gasteiger Group), Chinese Character Input File, Compressed Text, GE Industrial Systems CIMPLICITY Text Version HMI Screen, Alphacam Compiled Text (Planit), Online Course Text (Microsoft Corporation), Pretty Good Privacy (PGP) Ciphertext File (PGP Corporation), Windows Terminal Server INI Backup File (Microsoft Corporation), Visual Basic User Control Binary File, CTRAN/W DEFINE Compressed Data File (Geo-Slope International), Nokia PC Suite Backup Contact File (Nokia)\\
\end{tabular}

A package contains the following files:
\begin{itemize}
\item PACKAGE - package information (required, see \ref{sec:package information file})
\item FILES - list of files contained in the package (required for certification, see \ref{sec:files information file})
\item LICENSE files - the license files for the extension (required according to \ref{sec:package information file})
\item ... any files the developer might need below the virtual root "CONTAO" which refers to the contao installation root.
\item ... any files the developer might need below the virtual root "DATA" for the user data root.
\end{itemize}

\subsection[sec:package information file]{PACKAGE information file}

The XML format is described in the schema file extension.xsd\supref{lst:extension.xsd}.

As usage example serves extension.xml\supref{lst:extension.xml}.

\subsubsection[sec:dependencies conflicts suggestions]{Package relationships - dependencies, conflicts and preferations}

A relationship of an extension to another extension can be defined in form of a \techTerm{dependency}, a \techTerm{conflict} or a \techTerm{preferation}.
Each of these relationships can be towards any version of the mentioned extension or even be narrowed down to certain version numbers.

\subsubsection[sec:versioning of relationships]{Versioning of relationships}

As dependencies and conflicts are often regarding certain versions or a certain range of versions.

In the definition of an relationship, we may use the following attributes:
\begin{itemize}
\item \techTerm{match} which means an exact match on the version (X=Y).
\item \techTerm{lowest} which is implicit inclusive and tells the lower end of the range (X>=Y).
\item \techTerm{highest} which is implicit exclusive and tells the upper end of the range (X<Y).
\end{itemize}
The attributes \techTerm{lowest} and \techTerm{highest} do not have to be used in combination, therefore it is possible to define an open range.
When used in combination, they define a closed range.

The attributes \techTerm{lowest} and \techTerm{highest} must not be used in conjunction with the attribute \techTerm{match}. If used in conjunction with \techTerm{match}, the exact match takes precedence over the other attributes and therefore is always exact.

In the package information file, this might look like the following when defining a conflict for a certain db driver (in this case mysql).
\lstset{language=XML} 
\begin{lstlisting}[caption=XML example: conflict example mysql]
...
<ctx:conflict name="db:mysql">
    <ctx:version match="4"/>
    <ctx:version lowest="5" highest="5.1"/>
    <ctx:version lowest="5.2"/>
</ctx:conflict>
...
\end{lstlisting}

In above example definition we have a conflict for:
\begin{itemize}
\item all version 4 drivers.
\item all versions 5 to 5.1 ( 5>=X<5.1)
\item all versions from 5.2 on upwards (5.2>=X)
\end{itemize}
With above condition we define a conflict for all versions of MySQL except for version 5.1 (we could have done it more simple though).

\subsubsection[sec:handling of versions with extension component]{Handling of versions with extension component}

In general the extension component of a version is ignored as it makes common sense that an altered version of an extension is compatible to it's unaltered counterpart from an external point of view. This behaviour can be altered by defining the extension version in the relationship explicitely. If specified, the match for the version must be performed including the extension part.

\subsubsection[sec:special relationships]{Special relationships}

Package relationships can not only refer other packages but some other namespaces as well.
This is useful to restrict dependency on some php features, a certain database driver, some PEAR package or something else that has been defined in the specification.

Currently the following namespaces exist.\\
\begin{tabular}{ll}\\
contao         & The contao core.\\
php:[feature]  & Relates to a certain php feature. See php features in \ref{sec:package relationship namespaces php features}\\
pear:[package] & Relates to a certain pear package. See pear packages in \ref{sec:package relationship namespaces pear packages}.\\
db:type        & Relates to a certain database, \ie db:mysql. See database drivers in \ref{sec:package relationship namespaces database drivers}.\\
\end{tabular}

\subsubsection[sec:alternative dependencies]{Alternative dependencies}
Often it is required to have a certain feature available but multiple ways exist to get this feature available when depending on some PHP feature but to also provide a wrapper class as workaround if it is not available.

Instead of defining features within the repository client and therefore having to update the client all the time when a new feature has been defined, we provide the developer with the possibility to define alternative pathes by specifying groups of dependencies.

When defining groups of dependencies, the first positive match found is accepted and the group evaluation will be stopped. Note that the order in the XML has no meaning and therefore the attribute \techTerm{preferation} must be used to define which dependency is the most preferable.

\lstset{language=XML} 
\begin{lstlisting}[caption=XML example: alternative relationships]
...
<ctx:group preferation="php:soap nusoap">
<ctx:dependency name="php:soap">
    <ctx:description>I love soap, thats because i need this php extension.</ctx:description>
</ctx:dependency>
<ctx:dependency name="nusoap">
    <ctx:description>If you dont have soap, I pleased with an alternative implementation.</ctx:description>
</ctx:dependency>
</ctx:group>
...
\end{lstlisting}
This example describes, that it depends on the PHP SOAP extension but alternatively can also use the extension "nusoap" if the first is not available.

\subsubsection[sec:upgrade and downgrade compatibility]{Upgrade and downgrade compatibility}
In many cases upgrades and downgrades may cause problems. This is even more the case, when the data schema is changing.
That's the reason why the developer needs a way to define what potentional threat an upgrade or downgrade of a certain package implies.

This potentional threat can be categorized in the following sections:

\begin{tabular}{llp{0.55\textwidth}}\\
\textbf{Category} & \textbf{threat factor} & \textbf{Description}\\
lossless  & none & Up- and downgrade can be performed without data loss.\\
dropUnused  & almost none & Up- and downgrade will remove data that will not be in use anymore.\\
partial  & might be dangerous & Up- and downgrade might remove user data. Therefore manual intervention is suggested.\\
undefined & dragons live here! &This setting is implicit defined when the developer has not defined a setting at all. The result might be user data loss, in worst case scenario resulting in total data loss.\\
\end{tabular}

Whenever a developer has not defined a compatibility, the following automatic assumptions are made:\\
When upgrading, the compatibility is defaulting to "partial".
When downgrading, the compatibility is defaulting to "undefined".
The extension repository manager will then display a message where the up/downgrade compatibility is shown in addition to the description text of the compatibility.

As "undefined" is an implicit type, the developer can not add a description text. For "undefined" compatibility the extension repository client will provide two possibilities:
\begin{enumerate}
\item soft-update - upgrading without removing any user data (hence by updating the schema still some data might get lost).
\item hard-update - remove the old version including all user data and reinstall the other version then (possibly the worst attempt).
\end{enumerate}

A corresponding upgrade/downgrade compatibility might look like the following:
\lstset{language=XML} 
\begin{lstlisting}[caption=XML example: upgrade/downgrade compatibility]
...
<ctx:extension version="1.2.3">
    <ctx:upgrade version="1.1" userdata="lossless"/>
    <ctx:upgrade version="1.0" userdata="partial"/>
    <ctx:downgrade version="1.2" userdata="lossless"/>
    <ctx:downgrade version="1.1" userdata="dropUnused"/>
    <ctx:downgrade version="1.0" userdata="partial"/>
</ctx:extension>
...
\end{lstlisting}

\todo{translation}

Dabei beschreibt die Version immer die jüngste (Upgrade) oder älteste (Downgrade) Version, bis zu der diese Regel gilt. In dem Beispiel hat die Erweiterung die Version 1.2.3.
Die Upgrade Fähigkeit besagt:
Von allen Versionen 1.1 < 1.2.3 kann ohne Datenverlust aktualisiert werden.
Von allen Versionen 1.0 < 1.1 kann mit Nutzdatenverlust aktualisiert werden.
Die Downgrade Fähigkeit besagt:
Zu allen Versionen 1.2 < 1.2.3 kann ohne Datenverlust aktualisiert werden.
Zu allen Versionen 1.1 < 1.2 kann mit dem Verlust ungenutzer Daten aktualisiert werden.
Zu allen Versionen 1.0 < 1.1 kann mit Nutzdatenverlust aktualisiert werden.
Die Reihenfolge ist demnach egal, die Upgrade Definitionen werden zuerst nach Version absteigend sortiert und bilden dann immer zur nächsten Definition einen Versionsbereich, für den die Regel gilt.

\subsection[sec:files information file]{FILES - package contents information file}
\todo{translation}
Die Paketinhaltsinformation beinhaltet eine Liste aller Paketdateien (mit Ausnahme der FILES Datei selbst) mit samt Prüfsumme. Die Datei dient der Validierung der Installation. Sie ist erforderlich, wenn ein signiertes Paket erzeugt wird, um mit dem Extension Manager eine zuverlässige Integritätsprüfung zu gewährleisten. Bei einem unsigniertem Paket kann die Datei mitgeliefert werden, ist aber nicht erforderlich. Wird ein unsigniertes Paket installiert, erstellt der Paket Manager die Datei bei der Installation.
Beispieldatei

\lstset{language=XML} 
\begin{lstlisting}[caption=FILES example: package contents information file]
...
4543a0fe56f27277948469c15284ac25 CONTAO/system/modules/myExtension/Class1.php
a692ffb1c6026092d916872b2114941c CONTAO/system/modules/myExtension/Class2.php
2709656afab36a8f635206c6956cbf53 CONTAO/system/modules/myExtension/templates/class1\_template.tpl
3fced4c880f491dce98c116266361f1c CONTAO/system/modules/myExtension/templates/class2\_template.tpl
b52f2d57d10c4f7ee67a7eb9615d5d24 LGPL
2e7781b7b08ab489e08526f43adf5b0c PACKAGE
...
\end{lstlisting}





\section[sec:package relationship namespaces]{Package relationship namespaces}

\subsection[sec:package relationship namespaces php features]{PHP features}

\todo{this section is currently undocumented}

\subsection[sec:package relationship namespaces pear packages]{PEAR features}

\todo{this section is currently undocumented}

\subsection[sec:package relationship namespaces database drivers]{Database driver}

\todo{this section is currently undocumented}

\section[sec:namespaces]{Namespaces}

Contao 3 extensions should use the PHP 5.3 namespace feature. Contao 2 supports PHP prior 5.3 so this section only affects Contao 3 extensions.

\subsection[sec:namespace lock]{Namespace Lock}

Extensions can define a namespace lock, that means a namespace, that is "soft" locked for this extension.
This is a build in feature of the Extension Creator\supref{sec:extension builder} and should prevent multiple use of the same namespace.
Soft lock means, that it is not possible to exclusive lock a complete namespace, it is only possible to lock one namespace level.
For example if a developer lock the namespace \DisplFile{myextension\textbackslash} the namespace \DisplFile{myextension\textbackslash subnamespace\textbackslash} is not locked.
The developer have to lock every namespace he wants to use.

A developer who try to use a locked namespace get a warning that its locked from another extension.
Because this only affects the Extension Builder\supref{sec:extension builder}, it is not possible to absolutely prevent multiple use of the same namespace.


\chapter{Repository Stack}

In this chapter you find the whole extension repository stack defined.

\begin{landscape}
% I know, this is not really a state diagram but rather entities but hey.
\begin{figure}
\begin{emp}[Repository Stack](20, 20)

StateInfo.iRepositoryDistributionNetwork;
iRepositoryDistributionNetwork.foreColor := (.47, .94, .27);

NoteInfo.iUMLNote;
iUMLNote.foreColor := (0, .9, .9);


State.Portal("Repository Portal")();
ENote.NotePortal(iUMLNote)(
	"Public browsing interface",
	"- Public registration",
	"- Public profile management",
	"- Social networking");

State.RepManager("Repository Manager")();
ENote.NoteRepManager(iUMLNote)(
	"Authorization required management interface",
	"- User/Account management",
	"- Package builder",
	"- Package management",
	"- Administration");

State.RepMaster("Repository (Master)")();
ENote.NoteRepMaster(iUMLNote)("Package distribution network");

State.RepSlaveA("Repository (Slave1)")();
State.RepSlaveB("Repository (Slave2)")();
State.RepSlaveC("Repository (Slave3)")();

% this is ugly as hell but I was unable to position the notes more properly using
% same.left(NotePortal, NoteRepManager, NoteRepMaster);
% it always threw an error, FIX this.
leftToRight(95)(Portal, NotePortal);
leftToRight(90)(RepManager, NoteRepManager);
leftToRight(150)(RepMaster, NoteRepMaster);

% combine distribution network
leftToRight(20)(RepSlaveA, RepSlaveB);
topToBottom(20)(RepMaster, RepSlaveA);
EState.RepositoryDistributionNetwork(iRepositoryDistributionNetwork)("Package distribution network")(RepMaster, RepSlaveA, RepSlaveB);

% positioning of globals
topToBottom(20)(Portal, RepManager, RepositoryDistributionNetwork);

State.RepositoryStack("Repository Stack")(Portal, RepManager, RepositoryDistributionNetwork);

drawObjects(RepositoryStack, NotePortal, NoteRepManager, NoteRepMaster);

clink(transition)(RepMaster, RepSlaveA);
clink(transition)(RepMaster, RepSlaveB);

clink(transition)(Portal, RepManager);
clink(transition)(RepManager, RepMaster);


clink(dashedLink)(Portal, NotePortal);
clink(dashedLink)(RepManager, NoteRepManager);
clink(dashedLink)(RepMaster, NoteRepMaster);

\end{emp}
\caption{Repository Stack}
\end{figure}

\end{landscape}

\section[sec:repository]{Repository}
\todo{translation}

Das Repository ist der Paketspeicher, bietet Zugriff auf die Paketinformationen und verteilt die Pakete. Für den Zugriff und die Verwaltung der Pakete im Repository steht eine REST Schnittstelle zur Verfügung. Das Repository verfügt demnach nicht über ein eigenes Web-Interface.
Paketspeicher

Im Paketspeicher werden die Pakete abgelegt. Der Paketspeicher ist strukturiert aufgebaut, daraus ergeben sich folgende Paketpfade:
/<erster Buchstabe>/<name>/<name>-<version>.cext
z.B.
/a/Avisota/Avisota-1.5.1.cext
Die Pfade werden direkt auf die Installationswurzel gemapt.
REST Schnittstelle

Die REST Schnittstelle dient zum Abrufen von Informationen (ohne Authentifizierung) und zum Verwalten des Repositories (mit Authentifizierung).
Sie kann über dem Pfad /rest/ abgefragt werden.
Als Suffix wird das Abfrageformat angehängt, z.B. /json
Vor dem Abfrageformat werden die Parameter als Pfad-Komponente eingefügt, z.B. /rest/.../parameter/value/json
Schnittstellen

/rest/list/recent
Auflisten der neuesten Pakete.
Erlaubte Parameter: filter, offset, limit
/rest/list/categories
Auflisten der Kategorien in dem Repository.
Erlaubte Parameter: offset, limit
/rest/list/tags
Auflisten der Tags in dem Repository.
Erlaubte Parameter: offset, limit
/rest/search/keywords
Das Repository nach dem Schlüsselwörtern keywords durchsuchen.
Erlaubte Parameter: filter, offset, limit
/rest/put
Ein Paket in das Repository laden. Das Paket wird als POST Daten übertragen.
/rest/push
Ein paket in das Repository laden. Im Gegensatz zur put Methode wird in den POST Daten eine URL übertragen, vor der aus das Repository das Paket beziehen kann.
/rest/delete/name[/version]
Ein Paket aus dem Repository löschen. Das Paket wird über den Paketnamen identifiziert. Wenn die version angegeben wird, dann wird nur die angegebene Version gelöscht, ansonsten werden alle Versionen gelöscht.
/rest/settings/name/version
Liefert die Einstellungen (Kategorie, Veröffentlichung, Zugriff) einer Erweiterung im Repository.
/rest/configure/name/version
Die Einstellungen einer Erweiterung neu setzen. Als PUSH Daten werden die neuen Einstellungen übergeben.
Parameter

filter

Nach Kategorien filtern. Komma getrennte Liste.
z.B. /rest/list/rectent/filter/plugin/json
offset

Die ersten Treffer der Ergebnisliste überspringen.
z.B. /rest/list/recent/offset/10/json
limit

Das Ergebnis auf eine gewisse Anzahl Treffer begrenzen.
z.B. /rest/list/recent/limit/10/json
Ergebnisobjekt

Authentifizierung

Zum Zugriff auf die REST Schnittstelle gibt es 2 Authentifizierungsmethoden.
Authentifizierung mit Benutzername+Passwort
Authentifizierung mit AuthKey
Alle Zugriffe lassen sich zusätzlich per IP Filtern.
Setup Interface

Das Setup Interface ist das einzige Web-Interface, es dient zur Installation, Aktualisierung, Wartung und Verwaltung der Zugriffsrechte.


\section[sec:repository manager]{Repository Manager}

The repository manager is a web interface and the most intuitive form of managing a repository. The manager provides all basic functions of an repository and is part of the Portal\supref{sec:repository portal} but can also be used separately.

By combining the repository\supref{sec:repository package storage system} with the repository manager, one can build a simple to manage repository system. In fact this combination is in essence like the repository in version 2 on contao.org (but will supersede it in features and usability) from developer point of view. Extension packages will get built within the manager and published onto the repository via the REST API.

An installation that consists of this minimal installation is mainly of interest for agencies that want to run a private repository and do not intend to provide a public portal.

\section[sec:repository portal]{Repository Portal}

The portal is a public interface. It links external services to provide additional informations for extensions, that are available on the world wide web.

First it is a repository browser. Users can browse and search extensions and show details of them.
Secondary it is a information aggregator to display related informations.

A package contains links to other sources, see \ref{sec:package information file}.
\lstset{language=XML} 
\begin{lstlisting}[caption=XML example: links to other sources]
...
<ctx:link rel="homepage" href="http://www.contao.org">
    <ctx:description>Official project homepage.</ctx:description>
</ctx:link>
<ctx:link rel="manual" href="http://www.contao.org/documentation.html" language="en">
    <ctx:description>Official User-Manual.</ctx:description>
</ctx:link>
<ctx:link rel="help" href="http://en.contaowiki.org" language="en">
    <ctx:description>Inofficial Wiki.</ctx:description>
</ctx:link>
<ctx:link rel="video" href="http://www.youtube.com/watch?feature=player_embedded&amp;v=8HkQVmfjjYc">
    <ctx:description>Introduction into TypoLIGHT/Contao.</ctx:description>
</ctx:link>
<ctx:link rel="forum" href="http://www.contao-community.de" language="de">
    <ctx:description>German contao community.</ctx:description>
</ctx:link>
<ctx:link rel="tracker" href="http://dev.contao.org">
    <ctx:description>Official bug and issue tracker.</ctx:description>
</ctx:link>
<ctx:link rel="scm" href="http://svn.contao.org">
    <ctx:description>Official svn repository.</ctx:description>
</ctx:link>
...
\end{lstlisting}
The portal use this links and try to fetch related informations from these systems. I.e. this can be bugs and issues from a tracking system like Redmine or Github. Or screencast from YouTube. Or tutorials from a MediaWiki.

The Portal is not only an information aggregator. Not only the developer can link external informations to a extensions.
Other users can link ther own informations, like tutorials, screencasts or similar to an extension.
Theses links are stored in the portal, but not in the extension information. This is the social networking character of the Portal.

\subsection[sec:information kinds and system connectors]{Information kinds and system connectors}

The Portal use multiple connectors to get informations from other systems.
First we have to define the kind of informations, we are able to aggregate.

\begin{itemize}
\item Bugs and issues\\
The portal can display bugs and issues to show the current development status and upstream development.
\item Manuals and Tutorials - Wiki informations\\
One of the best ways to manage manuals and tutorials is to use a wiki. The portal can display related wiki articles.
\item Videos and Screencasts\\
Videos and screencasts are another kind of tutorials. The portal can display videos hosted on video streaming platforms like YouTube.
\item Source code\\
If an scm repository is configured, the portal can directly link to the source code.
\end{itemize}

These are just a few possible information kinds, the repository may use.

Next we have to define how the informations are aggregated.
A lot of systems like \hyperlink{http://www.mediawiki.org/wiki/API:Main_page/de}{MediaWiki}, \hyperlink{http://code.google.com/intl/de-DE/apis/youtube/overview.html}{YouTube} and \hyperlink{http://developer.github.com/v3/}{github} provide external APIs.
The Portal use connectors to this systems to get informations.
Because not all theses services are uniq, like MediaWiki the systems have to be configures dynamicaly.
All users are allowed to request adding a service, but only administrators are able to add them finaly into the portal.
Initialy there may be these connectors availeable:

\small
\begin{longtable}{|p{0.25\textwidth}|c|c|c|c}
\caption{Information system connectors and information kinds} \\
\hline
\label{tab:information system connectors and information kinds}
\textbf{System} & \textbf{Tracking} & \textbf{Wiki} & \textbf{Video} & \textbf{SCM}
\\ \hline
 MediaWiki & & X & &
\\ \hline
 Redmine & X & X & &
\\ \hline
 Github & X & X & & X
\\ \hline
 YouTube & & & X &
\\ \hline
 Vimeo & & & X &
\\ \hline
\end{longtable}
\normalsize


\label{sec:installation} 
\chapter{Installation}

This chapter describe how the Installation goes on and what programming basics are required.

\section[sec:er3 manager]{ER3 Manager}

The ER3 Manager (fka ER Client) is designed for Contao 2. It's a full replacement for the ER2 Client, that is currently in use.
It's a combined manager, that use extensions from the ER2 and the ER3. The ER3 Manager - or Extension Manager - have to manage all installed extension in an existing Contao 2 installation.

\section[sec:contao installation manager]{Contao Installation Manager}

The Contao Installation Manager (short CIM), is an Contao independend tool. It's designed for Contao 3 and should manage a complete Contao installation, including the installed Contao core and the installed extensions.
Summarized it's a combination of the current existing Install Tool and the ER3 Manager.

\paragraph{}
CIM should replace the Contao distribution package.
The idea behind CIM is, that a user only have to load a single PHP file, the CIM Installer.
The CIM Installer load all required CIM files.
Then the CIM provide a wizard to install a complete Contao.
Maybe with a basic template and basic extensions.

\paragraph{}
To make this possible, contao have to managed as a repository package itself.
This means, the CIM can manage the Contao core like any other extension.
In this case, there are no need for special behaviours to manage dependencies to the Contao Core.

\paragraph{}
But CIM is more than an installation tool.
CIM check the environment, if its compatible with Contao and provide solutions to make it compatible
For example activate the FTP file acces (fka SafeModeHack).
CIM manage the Contao system settings.
For example the database connection.
CIM simplified the installation process by automating the installation process.

\paragraph{}
For users, CIM is an absolute simplification of the installation process.
First download the CIM Installer (a small PHP file) and upload to the webspace.
Second call the CIM Installer, the CIM Installer download all required CIM files.
Third make shure, the environment is ready to install.
Fourth load a basic installation (Music Academy).
After CIM finish it work, the installation is ready to use.

\paragraph{}
For web workers/agencies additionaly it automating the installation process.
The automating process can be customized, to mind the custom preferences.


\section[sec:installation process]{Installation Process}

The Extension Manager (fka ER Client) use an installer class to install, upgrade, downgrade and uninstall extensions.
Every extension can use the basic installer, a modification of the basic installer or a complete own implementation.

An own implementation of the installer can be used to create individual installers.
For example the installer can fetch the extension files in the installation process.

TODO Problems with online-installation installer and signatures?!

\section[sec:installer interfaces and classes]{Installer interfaces and classes}

\begin{figure}[h]
\begin{emp}[Installer interfaces and classes](20, 20)
Class.information("InstallationInformation")()(
	"+getVersion() : Version",
	"+getFiles() : array",
	"+getInstalledPackage() : Package",
	"+getTargetPackage() : Package");

Class.installer("Installer")()(
	"+install($objInfo : InstallationInformation, $objVersion : Version) throws InstallerException : void",
	"+upgrade($objInfo : InstallationInformation, $objTargetVersion : Version) throws InstallerException : void",
	"+downgrade($objInfo : InstallationInformation, $objTargetVersion : Version) throws InstallerException : void",
	"+reinstall($objInfo : InstallationInformation) throws InstallerException : void",
	"+uninstall($objInfo : InstallationInformation) throws InstallerException : void",	
	"+clean($objInfo : InstallationInformation) throws InstallerException : void");

Class.basicInstaller("BasicInstaller")()(
	"+removeFiles($objInfo : InstallationInformation) throws InstallerException : void",
	"+removeContaoFiles($objInfo : InstallationInformation) throws InstallerException : void",
	"+removeDataFiles($objInfo : InstallationInformation) throws InstallerException : void",
	"+removeUnusedFiles($objInfo : InstallationInformation) throws InstallerException : void",
	"+removeUnusedContaoFiles($objInfo : InstallationInformation) throws InstallerException : void",
	"+removeUnusedDataFiles($objInfo : InstallationInformation) throws InstallerException : void",
	"+copyFiles($objInfo : InstallationInformation) throws InstallerException : void",
	"+copyContaoFiles($objInfo : InstallationInformation) throws InstallerException : void",
	"+copyUserFiles($objInfo : InstallationInformation) throws InstallerException : void",
	"+removeSchema($objInfo : InstallationInformation) throws InstallerException : void",
	"+safePatchSchema($objInfo : InstallationInformation) throws InstallerException : void",
	"+fullPatchSchema($objInfo : InstallationInformation) throws InstallerException : void");

topToBottom(20)(installer, basicInstaller);
topToBottom(20)(information, installer);

drawObjects(information, installer, basicInstaller);

clink(dashedLink)(information, installer);
clink(transition)(basicInstaller, installer);
\end{emp}
\caption{Installer interfaces and classes}
\end{figure}

\subsection[sec:installerinformation]{InstallationInformation}

The InstallationInformation object, contains all informations about an installation.

\begin{itemize}
\item Installed Version
\item Installed Files
\item Corresponding Schemas
\item The installed package file (if available)
\item The target installation package file (if a new version should be installed or up-/downgraded)
\end{itemize}

The InstallationInformation object is modified from the Installer.
The Installer register all it changes of the extension to the InstallationInformation object, including version changes, file list changes and schema changes.

\section[sec:installer]{Installer}

The installer have to provide methods to install, upgrade, downgrade and uninstall an extension.

\begin{figure}[h]
\begin{emp}[Installer Interface](20, 20)
Class.information("InstallationInformation")()(
	"+getVersion() : Version",
	"+getFiles() : array",
	"+getInstalledPackage() : Package",
	"+getTargetPackage() : Package");

Class.installer("Installer")()(
	"+install($objInfo : InstallationInformation, $objVersion : Version) throws InstallerException : void",
	"+upgrade($objInfo : InstallationInformation, $objTargetVersion : Version) throws InstallerException : void",
	"+downgrade($objInfo : InstallationInformation, $objTargetVersion : Version) throws InstallerException : void",
	"+reinstall($objInfo : InstallationInformation, $objTargetVersion : Version) throws InstallerException : void",
	"+uninstall($objInfo : InstallationInformation) throws InstallerException : void",	
	"+clean($objInfo : InstallationInformation) throws InstallerException : void");

topToBottom(20)(information, installer);

drawObjects(information, installer);

clink(dashedLink)(information, installer);
\end{emp}
\caption{Installer interface}
\end{figure}

\DisplCode{install(\$objInfo, \$objVersion)} Installs the extension. It gets the target version to install as argument to make it easer for developers to create multi version installers.

\vspace{.5cm}
\DisplCode{upgrade(\$objInfo, \$objTargetVersion)} Upgrades the extension from one version to another version.

\vspace{.5cm}
\DisplCode{downgrade(\$objInfo, \$objTargetVersion)} Downgrades the extension from one version to another version.

\vspace{.5cm}
\DisplCode{reinstall(\$objInfo, \$objTargetVersion)} Reinstall the extension with a specific version. The target version does not have to be the same, as the installed version. This method can be used to downgrade to a lower version, that is not supported by the \DisplCode{downgrade(..)} method.

\vspace{.5cm}
\DisplCode{uninstall(\$objInfo)} Uninstalls the extension. By definition, this will only uninstall the CONTAO files. All user DATA and Schema Information and data will be kept.

\vspace{.5cm}
\DisplCode{clean(\$objInfo)} Completely clean the installation, removes all files and schema data.

\vspace{.5cm}
The installation process success if there is no exception thrown.

\newpage
\section[sec:basic installer]{Basic Installer}

The Basic Installer is the default implementation of the installer interface.
If the developer does not define an installer for the extension, this installer is used.
In most cases, if the developer needs to add custom install, upgrade, downgrade or uninstall code, the developer create a supimplementation of this class.

\begin{figure}[h]
\begin{emp}[Basic Installer](20, 20)
Class.basicInstaller("BasicInstaller")()(
	"+removeFiles($objInfo : InstallationInformation) throws InstallerException : void",
	"+removeContaoFiles($objInfo : InstallationInformation) throws InstallerException : void",
	"+removeDataFiles($objInfo : InstallationInformation) throws InstallerException : void",
	"+removeUnusedFiles($objInfo : InstallationInformation) throws InstallerException : void",
	"+removeUnusedContaoFiles($objInfo : InstallationInformation) throws InstallerException : void",
	"+removeUnusedDataFiles($objInfo : InstallationInformation) throws InstallerException : void",
	"+copyFiles($objInfo : InstallationInformation) throws InstallerException : void",
	"+copyContaoFiles($objInfo : InstallationInformation) throws InstallerException : void",
	"+copyUserFiles($objInfo : InstallationInformation) throws InstallerException : void",
	"+removeSchema($objInfo : InstallationInformation) throws InstallerException : void",
	"+safePatchSchema($objInfo : InstallationInformation) throws InstallerException : void",
	"+fullPatchSchema($objInfo : InstallationInformation) throws InstallerException : void");

topToBottom(20)(installer, basicInstaller);

drawObjects(installer, basicInstaller);

clink(transition)(basicInstaller, installer);
\end{emp}
\caption{Basic Installer Implementation}
\end{figure}

% remove methods
\vspace{.5cm}
\DisplCode{removeFiles(\$objInfo)} Removes all CONTAO and DATA files from the installation.

\vspace{.5cm}
\DisplCode{removeContaoFiles(\$objInfo)} Removes all CONTAO files from the installation.

\vspace{.5cm}
\DisplCode{removeDataFiles(\$objInfo)} Removes all DATA files from the installation.

\vspace{.5cm}
\DisplCode{removeUnusedFiles(\$objInfo)} Removes unused CONTAO and DATA files from the installation.

\vspace{.5cm}
\DisplCode{removeUnusedContaoFiles(\$objInfo)} Removes unused CONTAO files from the installation.

\vspace{.5cm}
\DisplCode{removeUnusedDataFiles(\$objInfo)} Removes unused DATA files from the installation.

% copy files methods
\vspace{.5cm}
\DisplCode{copyContaoFiles(\$objInfo)} Copy the files from CONTAO package directory into the installation.

\vspace{.5cm}
\DisplCode{copyDataFiles(\$objInfo)} Copy the files from DATA package directory into the installation.

% schema methods
\vspace{.5cm}
\DisplCode{removeSchema(\$objInfo)} Remove the complete schema and all containing data from the database.

\vspace{.5cm}
\DisplCode{safePatchSchema(\$objInfo)} Add missing or modify existing schema, but does not delete any existing schema information and data.

\vspace{.5cm}
\DisplCode{fullPatchSchema(\$objInfo)} Add missing or modify existing schema, and delete no more existing schema information and data.

\vspace{.5cm}
The basic implementation may look like this. Please note, this is an incomplete, pseudo implementation!

\begin{lstlisting}[caption=Basic Installer Implementation Code]
class BasicInstaller implements Installer
{
	public function install(InstallationInformation \$objInfo, Version \$objVersion)
	{
		$this->copyContaoFiles($objInfo);
		$this->copyDataFiles($objInfo);
		$this->safePatchSchema($objInfo);
	}
	
	public function upgrade(InstallationInformation $objInfo, Version $objToVersion)
	{
		$this->removeUnusedFiles($objInfo);
		$this->safePatchSchema($objInfo);
		$this->copyFiles($objInfo);
	}
	
	public function downgrade(InstallationInformation $objInfo, Version $objToVersion)
	{
		$objVersionCompare = $objInfo->getVersion()->compare($objToVersion);
		
		// check if the versions differ in more than the build number
		if (!$objVersionCompare->isAnotherBuild())
		{
			throw new DowngradeException("Downgrade to another major, minor or revision is not supported.");
		}
		
		$this->removeUnusedFiles($objInfo);
		$this->safePatchSchema($objInfo);
		$this->copyFiles($objInfo);
	}
	
	public function reinstall(InstallationInformation $objInfo, Version $objTargetVersion)
	{
		$this->removeContaoFiles($objInfo);
		$this->safePatchSchema($objInfo);
		$this->copyFiles($objInfo);
	}
	
	public function uninstall(InstallationInformation $objInfo)
	{
		$this->removeContaoFiles($objInfo);
	}
	
	public function clean(InstallationInformation $objInfo)
	{
		$this->removeFiles($objInfo);
		$this->removeSchema($objInfo);
	}
	
	public function removeFiles(InstallationInformation $objInfo)
	{
		$this->removeContaoFiles($objInfo);
		$this->removeDataFiles($objInfo);
	}
	
	public function removeContaoFiles(InstallationInformation $objInfo)
	{
		// remove files from the CONTAO package directory.
		// unregister removed files from $objInfo
	}
	
	public function removeDataFiles(InstallationInformation $objInfo)
	{
		// remove files from the DATA package directory.
		// unregister removed files from $objInfo
	}
	
	public function removeUnusedFiles(InstallationInformation $objInfo)
	{
		$this->removeUnusedContaoFiles($objInfo);
		$this->removeUnusedDataFiles($objInfo);
	}
	
	public function removeUnusedContaoFiles(InstallationInformation $objInfo)
	{
		// remove files from the CONTAO package directory that are unused.
		// unregister removed files from $objInfo
	}
	
	public function removeUnusedDataFiles(InstallationInformation $objInfo)
	{
		// remove files from the DATA package directory that are unused.
		// unregister removed files from $objInfo
	}
	
	public function copyFiles(InstallationInformation $objInfo)
	{
		$this->copyContaoFiles($objInfo);
		$this->copyUserFiles($objInfo);
	}
	
	public function copyContaoFiles(InstallationInformation $objInfo)
	{
		// copy files from the CONTAO package into the installation.
		// register installed files into $objInfo
	}
	
	public function copyUserFiles(InstallationInformation $objInfo)
	{
		// copy files from the DATA package into the installation.
		// register installed files into $objInfo
	}
	
	public function removeSchema(InstallationInformation $objInfo)
	{
		// remove the complete schema
	}
	
	public function safePatchSchema(InstallationInformation $objInfo)
	{
		// add missing schema information and update existing
	}
	
	public function fullPatchSchema(InstallationInformation $objInfo)
	{
		// add missing schema information, update existing and remove no more existing informations.
	}
}
\end{lstlisting}

\section[sec:custom installer code]{Custom Installer Code}

If a developer want's to add custom code, he can make a subclass of the Basic Installer.

\begin{lstlisting}[caption=Custom Installer Code Example]
class CustomInstaller extends BasicInstaller
{
	public function upgrade(InstallationInformation $objInfo, Version $objToVersion)
	{
		// add custom pre-upgrade code here
		
		parent::upgrade($objInfo, $objToVersion);
		
		// add custom post-upgrade code here
	}	
}
\end{lstlisting}

\subsection[sec:basic installer]{Basic Installer}

The Basic Installer is the default implementation of the installer interface.
If the developer does not define an installer for the extension, this installer is used.
In most cases, if the developer needs to add custom install, upgrade, downgrade or uninstall code, the developer create a supimplementation of this class.

\begin{figure}[h]
\begin{emp}[Basic Installer](20, 20)
\end{emp}
\caption{Basic Installer Implementation}
\end{figure}

% remove methods
\vspace{.5cm}
\DisplCode{removeFiles(\$objInfo)} Removes all CONTAO and DATA files from the installation.

\vspace{.5cm}
\DisplCode{removeContaoFiles(\$objInfo)} Removes all CONTAO files from the installation.

\vspace{.5cm}
\DisplCode{removeDataFiles(\$objInfo)} Removes all DATA files from the installation.

\vspace{.5cm}
\DisplCode{removeUnusedFiles(\$objInfo)} Removes unused CONTAO and DATA files from the installation.

\vspace{.5cm}
\DisplCode{removeUnusedContaoFiles(\$objInfo)} Removes unused CONTAO files from the installation.

\vspace{.5cm}
\DisplCode{removeUnusedDataFiles(\$objInfo)} Removes unused DATA files from the installation.

% copy files methods
\vspace{.5cm}
\DisplCode{copyContaoFiles(\$objInfo)} Copy the files from CONTAO package directory into the installation.

\vspace{.5cm}
\DisplCode{copyDataFiles(\$objInfo)} Copy the files from DATA package directory into the installation.

% schema methods
\vspace{.5cm}
\DisplCode{removeSchema(\$objInfo)} Remove the complete schema and all containing data from the database.

\vspace{.5cm}
\DisplCode{safePatchSchema(\$objInfo)} Add missing or modify existing schema, but does not delete any existing schema information and data.

\vspace{.5cm}
\DisplCode{fullPatchSchema(\$objInfo)} Add missing or modify existing schema, and delete no more existing schema information and data.

\vspace{.5cm}
The basic implementation may look like this. Please note, this is an incomplete, pseudo implementation!

\begin{lstlisting}[caption=Basic Installer Implementation Code]
class BasicInstaller implements Installer
{
	public function install(InstallationInformation \$objInfo, Version \$objVersion)
	{
		$this->copyContaoFiles($objInfo);
		$this->copyDataFiles($objInfo);
		$this->safePatchSchema($objInfo);
	}
	
	public function upgrade(InstallationInformation $objInfo, Version $objToVersion)
	{
		$this->removeUnusedFiles($objInfo);
		$this->safePatchSchema($objInfo);
		$this->copyFiles($objInfo);
	}
	
	public function downgrade(InstallationInformation $objInfo, Version $objToVersion)
	{
		$objVersionCompare = $objInfo->getVersion()->compare($objToVersion);
		
		// check if the versions differ in more than the build number
		if (!$objVersionCompare->isAnotherBuild())
		{
			throw new DowngradeException("Downgrade to another major, minor or revision is not supported.");
		}
		
		$this->removeUnusedFiles($objInfo);
		$this->safePatchSchema($objInfo);
		$this->copyFiles($objInfo);
	}
	
	public function reinstall(InstallationInformation $objInfo, Version $objTargetVersion)
	{
		$this->removeContaoFiles($objInfo);
		$this->safePatchSchema($objInfo);
		$this->copyFiles($objInfo);
	}
	
	public function uninstall(InstallationInformation $objInfo)
	{
		$this->removeContaoFiles($objInfo);
	}
	
	public function clean(InstallationInformation $objInfo)
	{
		$this->removeFiles($objInfo);
		$this->removeSchema($objInfo);
	}
	
	public function removeFiles(InstallationInformation $objInfo)
	{
		$this->removeContaoFiles($objInfo);
		$this->removeDataFiles($objInfo);
	}
	
	public function removeContaoFiles(InstallationInformation $objInfo)
	{
		// remove files from the CONTAO package directory.
		// unregister removed files from $objInfo
	}
	
	public function removeDataFiles(InstallationInformation $objInfo)
	{
		// remove files from the DATA package directory.
		// unregister removed files from $objInfo
	}
	
	public function removeUnusedFiles(InstallationInformation $objInfo)
	{
		$this->removeUnusedContaoFiles($objInfo);
		$this->removeUnusedDataFiles($objInfo);
	}
	
	public function removeUnusedContaoFiles(InstallationInformation $objInfo)
	{
		// remove files from the CONTAO package directory that are unused.
		// unregister removed files from $objInfo
	}
	
	public function removeUnusedDataFiles(InstallationInformation $objInfo)
	{
		// remove files from the DATA package directory that are unused.
		// unregister removed files from $objInfo
	}
	
	public function copyFiles(InstallationInformation $objInfo)
	{
		$this->copyContaoFiles($objInfo);
		$this->copyUserFiles($objInfo);
	}
	
	public function copyContaoFiles(InstallationInformation $objInfo)
	{
		// copy files from the CONTAO package into the installation.
		// register installed files into $objInfo
	}
	
	public function copyUserFiles(InstallationInformation $objInfo)
	{
		// copy files from the DATA package into the installation.
		// register installed files into $objInfo
	}
	
	public function removeSchema(InstallationInformation $objInfo)
	{
		// remove the complete schema
	}
	
	public function safePatchSchema(InstallationInformation $objInfo)
	{
		// add missing schema information and update existing
	}
	
	public function fullPatchSchema(InstallationInformation $objInfo)
	{
		// add missing schema information, update existing and remove no more existing informations.
	}
}
\end{lstlisting}

\section[sec:custom installer code]{Custom Installer}

If a developer want's to add custom code, he can make a subclass of the BasicInstaller.

\begin{lstlisting}[caption=Custom Installer Code Example]
class CustomInstaller extends BasicInstaller
{
	public function upgrade(InstallationInformation $objInfo, Version $objToVersion)
	{
		// add custom pre-upgrade code here
		
		parent::upgrade($objInfo, $objToVersion);
		
		// add custom post-upgrade code here
	}	
}
\end{lstlisting}

Or the developer write his complete own installer. In this case the developer have to implement the Installer interface.


\label{sec:licensing} 
\chapter{Licensing}

There are two types of licenses, the multi license and the single license.
The multi license allow multiple installation of the extension.
The single license allow a single installation of the extension.

A license can be acquired through the "Repository Portal"\supref{sec:repository portal} or from the developer itself.
The "ER3 Manager"\supref{sec:er3 manager} allow to aquire a license directry through the "Repository Portal"\supref{sec:repository portal}.
To use licenses, the "ER3 Manager"\supref{sec:er3 manager} have to be authenticated against the "Repository Portal"\supref{sec:repository portal}.

The license is split into two parts, the license key and the license file.
The license key is a short unique identifier for the license file.
The license key have to be validated against the Repository (not Portal)

\section[sec:license key]{License key}

The license key is a sha256 hash, split into 8 parts with a length of 8 characters, concatenated with a minus character (-).

A license key may look like: \\
\texttt{e3b0c442-98fc1c14-9afbf4c8-996fb924-27ae41e4-649b934c-a495991b-7852b855}

\section[sec:license file]{License file}

The license file contains a lot of informations about the license, stored in INI format:
\begin{itemize}
\item Name of the extension [extension]
\item Lowest extension version [lowest version]
\item Highest extension version [highest version]
\item Name and email of license holder [license holder]
\item Name and email of licensor [licensor]
\item Unencrypted package checksum [checksum]
\item Type of license (single or multi, optional, default is multi) [type]
\end{itemize}
The names in the brackets [ ] are the ini fields.
The fields lowest version and highest version can be combined to a single field version.
This can be an absolut single version or a major or minor version.

\begin{lstlisting}[caption=License file]
[license]
extension=MyContaoExtension
lowest version=1.1
highest version=1
license holder=Tristan Lins <info@infinitysoft.de>
licensor=Christian Schiffler <info@cyberspectrum.de>
checksum=2c26b46b68ffc68ff99b453c1d30413413422d706483bfa0f98a5e886266e7ae
\end{lstlisting}

The license file is never stored in the contao extension. The license file can be optained from the "Repository"\supref{sec:repository} automaticaly by the "ER3 Manager"\supref{sec:er3 manager}. Alternatively the user can download the license file and use it every time, he want to install it (see Standalone license\supref{sec:standalone license}).

\section[sec:unencrypted package checksum]{Unencrypted package checksum}

The unencrypted package checksum (UPC) fulfilled two functions.
The first is to have an unique field in the license file, that cannot be recreated from a user.
The second is to validate the package integrity after decryption.

To calculate the UPC, the folowing steps are required.

Step 1: build a natural ordered list of all files, including the meta files of the package.

Step 2: build a sha256 hash about all concatenaded files, in the order of the list from step 1.

A sample implementation serves upc.php\supref{sec:unencrypted package checksum calculation}.

\section[sec:standalone license]{Standalone license}

A standalone license is used if the "ER3 Manager"\supref{sec:er3 manager} is unable to connect the repository or the repository does not contains the license. There is no different to a normal license, expecting that the license file is not obtained from the "Repository"\supref{sec:repository}. The user have to upload the license file from his computer every time, he want to install or update the extension. This user interaction is necessary because the license file cannot be stored in the contao installation to protect the package from steeling.

For this license usage the "License ipgrade"\supref{sec:license upgrade} is impossible.

\section[sec:encrypted packages]{Encrypted packages}

A licensed package get encrypted to prevent package stealing.
As encryption algorithm the advanced encryption standard (AES) is used.
If a package is encrypted, the \textbf{encrypted} flag have to be set to \textbf{true} and all encrypted files get the extension ".aes".
The encryption only affects files from the "CONTAO" and "DATA" directory.
All meta informations are stored unencrypted.
Encrypted packages requires signing of the packages.
The encryption key can be calculated from the license key and package checksum.
The license checksum (license key without minus) and package checksum are recalculated to bytestreams and then the bytestreams are XOR with each other. The result is the encryption key.
See upc.php\supref{sec:unencrypted package checksum calculation} for implementation example.

The usage of the licence checksum and package checksum as part of the encryption key forces every licensed package to be unique!

A sample implementation serves encrypt.php\supref{sec:package encryption implementation}.

\section[sec:single license]{Single license}

For a single installation there is used a third id to protect the package, the "installation id"\supref{sec:installation id}. For a single license the encryption key is calculated by XOR license checksum, package checksum and installation id. In the license file the "type" field have to set to "single". There are no other changes to the license, only the package use another key for encryption.

\subsection[sec:installation id]{Installation ID}

The installation id is unique for every installation. It is calculated from these factors:
\begin{itemize}
\item Contao Version \todo{Maybe its not the best idea to use Contao Version in the installation id?!}
\item Contao installation path (working directory)
\item Running domain
\end{itemize}

The installation id is a sha256 checksum over the three newline-concatenated values.

A sample implementation serves installationid.php\supref{sec:installation id calculation}.

\section[sec:license upgrade]{Licence upgrade}

As you read the sections above, there is a big problem with these licensing and encryption mechanism.
A license key can only be valid for a single package version.
This is because it is just a hash about the license file and the license file have to contain the checksum of the package.
As a result the "Repository Portal"\supref{sec:repository portal} have to store the information for which version the license is generated.
This is the basic for the license upgrade.

The licensed (encrypted) packages are automatic generated through the "Extension Builder"\supref{sec:extension builder}.
Because this, its possible to automatic generate a new license and package.
If a user try to update a licensed extension and does not own a license for the current version, but for a prior version.
Then the "Repository Portal"\supref{sec:repository portal} automaticale upgrade the license.
In the license file contains all necessary informations for this.
If an upgrade is possible, this means the license includes the current version, a new license file and license key is generated.
The license key is pushed to and stored in the contao installation.
After that, the "ER3 Manager"\supref{sec:er3 manager} can download and upgrade to the new package version.

If the "Extension Builder"\supref{sec:extension builder} is not used by the developer, the developer have to upload an unencrypted package into the repository. The unencrypted package will never be published, but its used to generate the licensed package for every license holder.

% UML support
\end{empfile}

% Seitennummerierung -----------------------------------------------------------
%   Nach dem Hauptteil werden die Seiten in großen römischen Ziffern 
%   nummeriert.
% ------------------------------------------------------------------------------
\pagenumbering{Roman}

% Anhang -----------------------------------------------------------------------
%   Die Inhalte des Anhangs werden analog zu den Kapiteln inkludiert.
%   Dies geschieht in der Datei "Appendix.tex".
% ------------------------------------------------------------------------------
\begin{appendix}
    \clearpage
    \chapter{Appendix}
    \label{sec:Appendix}
    \pagenumbering{roman}
    % Rand der Aufzählungen in Tabellen anpassen
    \setdefaultleftmargin{0em}{}{}{}{}{}
    \section[sec:package information file schema]{PACKAGE information file schema}
\xmllisting{lst:extension.xsd}{PACKAGE information file schema}{../extension.xsd}

\section[sec:package information file example]{PACKAGE information file example}
\xmllisting{lst:extension.xml}{PACKAGE information file example}{../extension.xml}

\section[sec:unencrypted package checksum calculation]{Unencrypted package checksum calculation}
\phplisting{lst:upc.php}{Unencrypted package checksum calculation}{../../src/upc.php}

\section[sec:package encryption implementation]{Package encryption}
\phplisting{lst:encrypt.php}{Package encryption}{../../src/encrypt.php}

\section[sec:installation id calculation]{Installation id calculation}
\phplisting{lst:installationid.php}{Installation id calculation}{../../src/installationid.php}

\end{appendix}

% Abkürzungsverzeichnis --------------------------------------------------------
\input{content/Glossar}
% für korrekte Überschrift in der Kopfzeile
\clearpage\markboth{\nomname}{\nomname} 
\printnomenclature
\label{sec:Glossar}

\listoffigures % Abbildungsverzeichnis
\listoftables % Tabellenverzeichnis
\renewcommand{\lstlistlistingname}{List of Listings}
\lstlistoflistings % Listings-Verzeichnis

% Literaturverzeichnis ---------------------------------------------------------
%   Das Literaturverzeichnis wird aus der BibTeX-Datenbank "Bibliographie.bib"
%   erstellt.
% ------------------------------------------------------------------------------
\bibliography{Bibliographie} % Aufruf: bibtex RFCCC1
\bibliographystyle{natdin} % DIN-Stil des Literaturverzeichnisses

% Index ------------------------------------------------------------------------
%   Zum Erstellen eines Index, die folgende Zeile auskommentieren.
% ------------------------------------------------------------------------------
\printindex
\printglossaries

\end{document}
